\documentclass[a4paper,12pt]{article}

\usepackage[a4paper,margin=2.3cm,headheight=16pt]{geometry}

\usepackage[
backend=biber,
style=ieee
]{biblatex}

\addbibresource{references.bib}

\usepackage{xepersian}
\usepackage{fancyhdr}
\usepackage{array}
\usepackage{setspace}

%---------------- Fonts ----------------%
\settextfont{XB Niloofar}
%---------------- Paragraph Settings ----------------%
\setlength{\parindent}{0pt}
\setlength{\parskip}{0.45em}
\linespread{1.08}

%---------------- Header ----------------%
\pagestyle{fancy}
\fancyhf{}
\fancyhead[R]{\textbf{نام درس:} رایانش ابری}
\fancyhead[L]{\textbf{استاد:} دکتر امیر کلباسی}
\renewcommand{\headrulewidth}{0.4pt}
\renewcommand{\footrulewidth}{0pt}

\begin{document}
	
	\begin{center}
		{\Large \textbf{گزارش پایانی پروژه}}
		
		\vspace{0.45cm}
		
		{\large \textbf{موضوع}}
		
		\vspace{0.2cm}
		
		\textbf{طراحی و ارزیابی رویکردی برای بهبود زمان‌بندی وظایف یادگیری ماشین در Kubernetes}
		
		\vspace{0.2cm}
		
		{\small سامانهٔ پیاده‌سازی‌شده: \lr{KubeML-Aware}}
		
		\vspace{0.6cm}
		
		\begin{center}
			\begin{minipage}{0.42\textwidth}
				\centering
				\textbf{امیر میرزائی موحد}\\[0.15cm]
				شماره دانشجویی: 404125059
			\end{minipage}
			\hfill
			\begin{minipage}{0.42\textwidth}
				\centering
				\textbf{مجید بافنده}\\[0.15cm]
				شماره دانشجویی: 404131043
			\end{minipage}
		\end{center}
	\end{center}
	
	\vspace{0.35cm}
	
	{\large\textbf{۱. خلاصهٔ پروژه}}
	
	در این پروژه یک لایهٔ \emph{پذیرش آگاه از رتبه} برای بارهای کاری آموزش یادگیری ماشین در \lr{Kubernetes} با نام \lr{KubeML-Aware} طراحی و پیاده‌سازی شد. برخلاف رویکردهایی که زمان‌بند پیش‌فرض را جایگزین می‌کنند، این سامانه تنها \emph{زمان واجد شرایط ‌شدن} پادها\footnote{\lr{Pod}: واحد اجرایی در Kubernetes} برای زمان‌بندی را کنترل می‌کند و وظایف فیلترکردن گره‌ها، امتیازدهی و اتصال\footnote{\lr{Pod}} را همچنان به زمان‌بند پیش‌فرض واگذار می‌کند. سازوکار کار به این صورت است که همهٔ پادهای یک دسته\footnote{\lr{batch}} پشت یک دروازهٔ زمان‌بندی پایدار نگه داشته می‌شوند؛ پس از کامل‌شدن دسته، یک کنترلر شش ویژگی هر وظیفه را اعتبارسنجی می‌کند، یک رتبهٔ قطعی\footnote{\lr{deterministic}} محاسبه می‌کند و دروازه‌ها را به ترتیب رتبه و بر اساس بازخورد منابع برمی‌دارد. این طراحی امکان بهبود ترتیب اجرای وظایف را بدون از دست دادن ایمنی و پایداری زمان‌بند اصلی فراهم می‌کند.
	
	{\large\textbf{۲. دستاوردهای اصلی}}
	
	مهم‌ترین خروجی‌های عملی پروژه عبارت‌اند از: الگوریتم رتبه‌بندی شش‌ ویژگی با نرمال‌سازی درون‌دسته‌ای و مدیریت رتبه‌های یکسان؛ سازوکار \lr{FastPath} که با استفاده از تله‌متری \footnote{\lr{telemetry}} تازهٔ منابع، فاصله‌های خالی صف زمان‌بند را حذف می‌کند و در صورت کهنه‌بودن داده‌ها به‌صورت ایمن (\lr{fail-closed}) عمل می‌کند؛ یک \lr{tail-balancer} برای بازچینش وظایف کم‌اهمیت‌تر؛ بستهٔ استقرار \lr{Helm} به همراه تنظیمات \lr{RBAC}؛ و در نهایت اجرای دو دور آزمایش واقعی روی خوشهٔ \lr{K3s}. در ارزیابی آزمایشی اولیه (دوازده وظیفه، پنج تکرار جفتی)، میانگین زمان تکمیل وظیفه\footnote{\lr{JCT}} بهبود ۷٫۳۸ درصدی داشت. در ارزیابی نهایی و در مقیاس بزرگ‌ یک دستهٔ ۲۰۰ وظیفه‌ای، میانگین \lr{JCT} حدود \textbf{۳۰ درصد}، صدک ۹۵ زمان تکمیل (\lr{p95 JCT}) حدود \textbf{۱۳٫۷ درصد} و زمان کل اتمام دسته (\lr{makespan}) حدود \textbf{۱۲٫۸ درصد} نسبت به زمان‌بند پیش‌فرض \lr{Kubernetes} بهبود یافتند.
	
	\vspace{0.2cm}
	
	{\large\textbf{۳. فعالیت‌های انجام‌شده}}
	
	با توجه به دو نفره بودن پروژه، حجم کار به‌گونه‌ای انتخاب شد که فراتر از یک پروژهٔ تک‌نفره باشد: علاوه بر طراحی الگوریتم و نگارش مقاله، یک سامانهٔ کامل و آمادهٔ استقرار در \lr{Kubernetes} همراه با مجموعهٔ گسترده‌ای از آزمون‌ها پیاده‌سازی و بر روی یک سرور ارزیابی شد. فعالیت‌های پروژه را می‌توان در چهار حوزهٔ اصلی زیر دسته‌بندی کرد:
	
	\textbf{الف) پیاده‌سازی و توسعهٔ سامانه.}
	در این بخش، هستهٔ کنترلر دروازه‌ای (\lr{gate controller})، الگوریتم رتبه‌بندی (\lr{rank}), منطق اجرا و ثبت رویدادها (\lr{execution}, \lr{records})، سازوکار جمع‌آوری دسته (\lr{burst})، مسیر سریع (\lr{fast\_path})، تنظیم آهنگ تطبیقی (\lr{pacing}) و مدل کار مصنوعی آموزش (\lr{work\_model}) توسعه یافتند. همچنین بستهٔ \lr{Helm}، منابع \lr{RBAC} و مانیفست‌های \lr{Knative} برای استقرار امن آماده شدند.
	
	\textbf{ب) آزمایش و جمع‌آوری داده.}
	مجموعه‌ای از بارهای کاری متنوع (کوچک، متوسط، بزرگ، ناهمگون و ابلیشن\footnote{\lr{ablation}}) طراحی شد و پروتکل آزمایش جفتی\footnote{\lr{paired}} توسط \lr{K3s} با پنج تکرار و دوازده وظیفه در هر بازو\footnote{\lr{arm}} اجرا گردید. کد جمع‌آوری متریک‌ها (\lr{metrics\_collector})، تله‌متری و ثبت رویدادهای سطح‌برنامه برای اندازه‌گیری دقیق \lr{JCT} و \lr{makespan} توسعه داده شد. در گام بعد، برای افزایش توان آماری و بررسی رفتار سامانه در مقیاس واقعی‌تر، یک دور دوم آزمایش با ۲۰۰ وظیفهٔ یادگیری ماشین ناهمگون که با شش ویژگی رتبه‌بندی (مدت‌زمان تخمینی، نرخ کاهش خطا، بُعد ماتریس، حجم به‌روزرسانی گرادیان، فاصلهٔ نقطهٔ بازرسی و تعداد افراز مدل) توصیف می‌شدند به‌صورت یک دستهٔ واحد روی \lr{K3s} اجرا و در چهار پیکربندی مقایسه شد: زمان‌بند پیش‌فرض \lr{Kubernetes}، رتبه‌بندی تک‌ویژگی بر مبنای مدت‌زمان (\lr{duration-only})، الگوریتم کامل شش‌ویژگی \lr{KubeML-Aware}، و یک کنترل ابلیشن با رتبه‌بندی معکوس‌شده (برای بررسی جهت اثر).
	
	\textbf{ج) تحلیل نتایج و نگارش مقاله.}
	داده‌های خام آزمایش با اسکریپت‌های تحلیل آماری پردازش شد. سپس مقالهٔ نهایی به سبک \lr{IEEE} به همراه شکل‌های معماری، بافت \lr{Kubernetes}، نمودار نتایج و خط‌زمانی اجرا تدوین و در چند مرحله بازبینی شد.
	
	\textbf{د) آزمون و رفع اشکال.}
	یک مجموعهٔ جامع از آزمون‌های واحد و یکپارچگی برای همهٔ بخش‌ها (رتبه‌بندی، کنترلرها، اجرا، مسیر سریع، تله‌متری، تولید بار کاری و اعتبارسنجی مانیفست‌ها) نوشته شد. در این مرحله باگ‌های مربوط به تجزیهٔ لاگ‌ها شناسایی و برطرف شد و صحت مانیفست‌های \lr{Kubernetes} پیش از استقرار اعتبارسنجی گردید.
	
	\vspace{0.2cm}
	
	{\large\textbf{۴. نقش و سهم هر یک از اعضا}}
	
	مطابق خواستهٔ درس، نقش هر عضو در ادامه به‌صورت دقیق و صریح مشخص شده است. نکتهٔ اصلی این است که هر دو عضو در تمام بخش‌های پروژه شامل نگارش مقاله، کدنویسی، اجرای آزمایش‌ها و رفع اشکال به‌صورت فعال مشارکت داشته‌اند و هیچ بخشی به‌طور کامل بر عهدهٔ تنها یک نفر نبوده است. تصمیم‌های اصلی طراحی، ساختار مقاله و تحلیل نتایج به‌صورت مشترک اتخاذ شده است. با این حال، برای شفافیت بیشتر، \emph{مسئولیت اصلی/راهبری} هر حوزه به شکل زیر تقسیم شده است:
	اعضای گروه به‌صورت مشترک در تمامی مراحل پروژه مشارکت داشتند. هر دو عضو در طراحی و پیاده‌سازی هستهٔ کنترلر دروازه‌ای، الگوریتم رتبه‌بندی شش‌ویژگی، سازوکار \lr{FastPath}، طراحی بارهای کاری، اجرای آزمایش‌ها روی \lr{K3s}، جمع‌آوری و تحلیل داده‌ها، توسعه و اجرای آزمون‌ها، رفع اشکال و نگارش و بازبینی مقاله نقش فعال داشتند. تصمیم‌های فنی، طراحی‌های اصلی و تحلیل نتایج نیز به‌صورت مشترک انجام شد و هیچ‌یک از بخش‌های پروژه به‌طور مستقل به یکی از اعضا اختصاص نداشت.

	
	\vspace{0.2cm}
	
	{\large\textbf{۵. نتایج کلیدی}}
	
	{\large\textbf{۵.۱ ارزیابی پایلوت (دوازده وظیفه، پنج تکرار)}}
	
	دور نخست ارزیابی روی یک گرهٔ چهارهسته‌ای اشتراکی \lr{K3s} و با پنج تکرار جفتی انجام شد. نتایج در جدول زیر خلاصه شده‌اند. میانگین زمان تکمیل وظیفه در تمام پنج جفت بهبود یافت و بهبود میانگین جفتی برابر ۷٫۳۸ درصد (با بازهٔ اطمینان اکتشافی ۹۵ درصد برابر [۲٫۵۷، ۱۲٫۰۵]) به‌دست آمد. زمان کل اتمام دسته (\lr{makespan}) نیز در هر پنج جفت در همان جهت کاهش یافت (۲٫۴۷ درصد)، اما چون بازهٔ اطمینان آن شامل صفر است، در آن مرحله به‌عنوان یک روند جهت‌دار و نه یک اثر قطعی گزارش شد.
	
	\begin{center}
		\setlength{\arrayrulewidth}{0.35pt}
		\renewcommand{\arraystretch}{1.3}
		\begin{tabular}{|>{\raggedright\arraybackslash}p{3.6cm}|c|c|c|}
			\hline
			\textbf{معیار} & \textbf{\lr{FIFO} پایه} & \textbf{\lr{KubeML-Aware}} & \textbf{بهبود} \\
			\hline
			میانگین \lr{JCT} & ۳۱٫۰۵ ثانیه & ۲۸٫۷۶ ثانیه & ‎+۷٫۳۸٪‎ \\
			\hline
			\lr{Makespan} & ۵۵٫۳۷ ثانیه & ۵۴٫۰۱ ثانیه & ‎+۲٫۴۷٪‎ \\
			\hline
		\end{tabular}
	\end{center}
	
	{\large\textbf{۵.۲ ارزیابی نهایی در مقیاس بزرگ (۲۰۰ مرحله)}}
	
	از آنجا که نتایج پایلوت امیدوارکننده ولی از نظر آماری محدود بودند (به‌ویژه برای \lr{makespan})، دور دوم آزمایش با هدف افزایش توان آماری و سنجش پایداری اثر در مقیاس بزرگ‌تر طراحی شد. در این دور، یک دستهٔ ۲۰۰ وظیفهٔ یادگیری ماشین ناهمگون روی همان سرویس \lr{K3s} اجرا و در چهار پیکربندی مقایسه گردید: زمان‌بند پیش‌فرض \lr{Kubernetes} (خط پایه)، رتبه‌بندی تک‌ویژگی \lr{duration-only}، الگوریتم کامل شش‌ویژگی \lr{KubeML-Aware}، و یک کنترل ابلیشن با رتبهٔ معکوس‌شده. نتایج نسبت به خط پایه در جدول زیر آمده است:
	
	\begin{center}
		\setlength{\arrayrulewidth}{0.35pt}
		\renewcommand{\arraystretch}{1.3}
		\begin{tabular}{|>{\raggedright\arraybackslash}p{3.4cm}|c|c|c|}
			\hline
			\textbf{پیکربندی} & \textbf{میانگین \lr{JCT}} & \textbf{\lr{p95 JCT}} & \textbf{\lr{Makespan}} \\
			\hline
			\lr{Kubernetes} پیش‌فرض (خط پایه) & ۹۹۰٫۸ ثانیه & ۱۲۹۶٫۶ ثانیه & ۱۳۰۴٫۵ ثانیه \\
			\hline
			رتبه‌بندی \lr{duration-only} & ۷۳۷٫۹ ثانیه (‎+۲۵٫۵٪‎) & ۱۱۵۳٫۲ ثانیه (‎+۱۱٫۱٪‎) & ۱۱۶۰٫۸ ثانیه (‎+۱۱٫۰٪‎) \\
			\hline
			\textbf{\lr{KubeML-Aware} (کامل)} & \textbf{۶۹۳٫۶ ثانیه (‎+۳۰٫۰٪‎)} & \textbf{۱۱۱۸٫۶ ثانیه (‎+۱۳٫۷٪‎)} & \textbf{۱۱۳۷٫۶ ثانیه (‎+۱۲٫۸٪‎)} \\
			\hline
			ابلیشن رتبهٔ معکوس & ۱۲۳۵٫۵ ثانیه (‎-۲۴٫۷٪‎) & ۱۴۴۳٫۴ ثانیه (‎-۱۱٫۳٪‎) & ۱۴۴۷٫۱ ثانیه (‎-۱۰٫۹٪‎) \\
			\hline
		\end{tabular}
	\end{center}
	
	در این مقیاس، الگوریتم کامل \lr{KubeML-Aware} در هر سه معیار بهترین عملکرد را داشت: میانگین \lr{JCT} حدود \textbf{۳۰ درصد}، \lr{p95 JCT} حدود \textbf{۱۳٫۷ درصد} و \lr{makespan} حدود \textbf{۱۲٫۸ درصد} نسبت به زمان‌بند پیش‌فرض بهبود یافتند. نکتهٔ قابل‌توجه این است که رتبه‌بندی شش‌ویژگی به‌طور پیوسته از رتبه‌بندی تک‌ویژگی \lr{duration-only} بهتر عمل کرد، که نشان می‌دهد ویژگی‌های اضافی (نرخ کاهش خطا، بُعد ماتریس، حجم به‌روزرسانی گرادیان، فاصلهٔ نقطهٔ بازرسی و تعداد افراز مدل) اطلاعات مفیدی فراتر از مدت‌زمان تخمینی به الگوریتم می‌افزایند. کنترل ابلیشن با رتبهٔ معکوس‌شده، در هر سه معیار عملکردی به‌طور محسوسی بدتر از خط پایه داشت (کاهش عملکرد بین ۱۱ تا ۲۵ درصد)؛ این الگو یک شاهد مهم است که بهبود مشاهده‌شده ناشی از \emph{جهت درست} سیگنال رتبه‌بندی است و نه صرفاً اثر جانبی مکانیزم دروازه‌بندی یا سربار زمان‌بندی.
	
	\vspace{0.2cm}
	
	{\large\textbf{۶. تحلیل تکمیلی و اعتبارسنجی}}
	
	برای اطمینان از قابل‌اتکابودن نتایج مقیاس بزرگ، سه بررسی تکمیلی انجام شد. نخست، توزیع تجمعی تجربی (\lr{ECDF}) زمان تکمیل وظایف در هر چهار پیکربندی رسم و مقایسه شد؛ منحنی \lr{KubeML-Aware} در بیشتر بازهٔ صدک‌ها به‌طور پیوسته سمت چپ منحنی خط پایه قرار دارد که نشان‌دهندهٔ بهبود یکنواخت و نه فقط جابه‌جایی میانگین است. دوم، منحنی تکمیل تجمعی وظایف بر حسب زمان از شروع دسته بررسی شد که نشان داد \lr{KubeML-Aware} با وجود شروع کندتر در وظایف نخست (به دلیل انتظار برای تکمیل‌شدن دسته پشت دروازهٔ زمان‌بندی)، به‌سرعت روند سریع‌تری در تکمیل وظایف باقی‌مانده پیدا می‌کند و در پایان دستهٔ ۲۰۰ وظیفه‌ای زودتر از خط پایه به اتمام می‌رسد. سوم، سازگاری جهت اثر بین دو دور آزمایش (پایلوت و مقیاس بزرگ) و همچنین بین سه معیار مختلف (میانگین، صدک ۹۵ و \lr{makespan}) بررسی شد؛ در هر دو دور و هر سه معیار، پیکربندی \lr{KubeML-Aware} به‌طور پیوسته بهتر از خط پایه و بهتر از نسخهٔ \lr{duration-only} بود، که اطمینان بیشتری نسبت به پایداری اثر مشاهده‌شده فراهم می‌کند.
	
	\vspace{0.2cm}
	
	{\large\textbf{۷.  جمع‌بندی}}
	
	در مجموع، این پروژه نشان داد که می‌توان یک سیاست پذیرش آگاه از ویژگی‌های یادگیری ماشین را بدون جایگزینی موتور مکان‌یابی \lr{Kubernetes} افزود و به بهبود قابل‌مشاهده و پایدار در میانگین زمان تکمیل وظایف، صدک ۹۵ و زمان کل اتمام دسته رسید. نتایج پایلوت (بهبود ۷٫۳۸ درصدی در مقیاس کوچک) و نتایج مقیاس بزرگ (بهبود حدود ۳۰ درصدی در میانگین \lr{JCT}) هم‌جهت هستند و اعتماد بیشتری به اثر واقعی الگوریتم رتبه‌بندی می‌دهند. گام بعدی، اجرای آزمایش روی خوشهٔ اختصاصی (به‌جای گرهٔ اشتراکی) با تکرارهای بیشتر و بارهای کاری واقعی چارچوب‌های یادگیری ماشین (به‌جای مدل‌سازی‌شده) است تا این نتایج در شرایط تولید نیز تأیید شوند.
	
\end{document}